\appendixchapter{Problem Indicators}
\label{appendix:problem_indicators}

\begin{table}[H]
  \caption{Quantitative Evidence Supporting the Need for Alternative Input Approaches}
  \label{tbl:problem_indicators}
  \begin{tabularx}{\textwidth}{|Y|Y|Y|}
    \hline
    \textbf{Dimension} & \textbf{Evidence} & \textbf{Interpretation} \\
    \hline
    Population burden & Work-related MSD prevalence among computer users: 33.8\%--95.3\%~\cite{demissie2024wrmsd} & The risk profile is broad and persistent across user groups. \\
    \hline
    Global trajectory & MSD cases increased 123.4\% (1990--2020), with continued growth projected~\cite{gbd2023othermsd} & The problem scale is increasing, not stabilising. \\
    \hline
    Biomechanical exposure & Mouse use $>$50\% of work time increases hand/wrist symptom odds ratio (OR) to 1.68~\cite{andersen2003mouse} & High-exposure workflows amplify cumulative upper-limb risk. \\
    \hline
    Clinical signal & CTS prevalence 7.8\% in pooled occupational data~\cite{dale2013cts} & A major musculoskeletal outcome is measurable in occupational populations. \\
    \hline
    Accessibility barrier & 85\% of arthritis participants report major computer-use difficulties~\cite{baker2009arthritis} & Classical input assumptions do not generalise across motor conditions. \\
    \hline
    Interaction performance gap & Fitts' law throughput: mouse 4.73~bits/second (bps) vs.\ gyrosensor 2.85~bps~\cite{ramcharitar2017fitts} & Spatial alternatives are viable but require algorithmic refinement to approach parity. \\
    \hline
  \end{tabularx}
\end{table}